\section{Conclusión}

\begin{itemize}
    \item La implementación del sistema Odoo ha permitido a PERÚ FLORES TRAVEL AGENCY TOURS OPERADOR E.I.R.L. optimizar sus procesos comerciales y operativos, logrando una integración efectiva que elimina los "silos de información" antes dispersos en Excel, WhatsApp y notas manuales. Esta centralización garantiza ahora una trazabilidad completa de cada pasajero, desde el primer contacto hasta la post-venta.
    \item El sistema ha contribuido a estandarizar la oferta comercial mediante la digitalización del catálogo de paquetes turísticos (Cusco, Arequipa, Puno). Esto ha permitido una reducción drástica en los tiempos de respuesta, logrando generar cotizaciones profesionales en PDF en menos de 15 minutos, comparado con los 45-60 minutos que tomaba el proceso manual.
    \item La digitalización de los procesos ha fortalecido la calidad del servicio al reducir la tasa de error operativo a menos del 1\%. Asimismo, la implementación del módulo de Encuestas permite a la agencia capturar la satisfacción del cliente de manera sistemática, generando información valiosa para fortalecer su reputación digital y el "boca a boca" en su sitio web.
    \item Si bien la adopción del sistema implicó un desafío cultural ---pasando de una "flexibilidad manual" a una "estandarización digital"---, este proceso fue superado con éxito mediante capacitaciones prácticas dirigidas a la Sra. Silveria y su equipo, asegurando el dominio de la plataforma tanto en escritorio como en dispositivos móviles.
    \item Finalmente, el proyecto ha sentado las bases para la escalabilidad del negocio en el competitivo mercado turístico de Cusco. Con una infraestructura digital moderna, la agencia se encuentra ahora en una posición favorable para tomar decisiones estratégicas basadas en datos reales y para implementar futuras mejoras como pagos en línea y automatización avanzada de marketing.
\end{itemize}
